\begin{pastexam}[label={px:ncku-112-cce-ds-05}]{成大電通所 112 學年度 第 5 題}
（問答）

(10\%) Landscapes in films are often computer generated. Ever wondered how they do it? Here is an
algorithmic drawing example. Draw enough repeating patterns until you can tell what it is. Start by
drawing a single straight vertical line as below

\begin{center}
\begin{tikzpicture}
\draw[thick] (0,0)--(0,1.6);
\end{tikzpicture}
\end{center}

Then execute Draw( ) from this vertical line until you can show how the algorithmic drawing works.

\begin{center}
\fbox{\begin{minipage}{0.9\textwidth}\sffamily
Draw( )\\
\{\\[2pt]
{\footnotesize
\hspace*{1.5em}Draw 3 shorter lines at an angle in the top two-thirds of the line on its left side.\\[2pt]
\hspace*{1.5em}Draw 3 shorter lines at an angle in the top two-thirds of the line on its right side.\\[2pt]
\hspace*{1.5em}Choose a new existing line and Draw( ) from that line again\\[2pt]}
\}
\end{minipage}}
\end{center}
\end{pastexam}

\begin{pastsoln}
本題和成大智慧資安學程同年的第~5~題相同（三個系所共用同一份考卷），完整的詳解見歷屆試題~\ref{px:ncku-112-isec-ds-05}，這裡只列答案。

答案：畫出來的是一棵樹（像松樹、杉樹那樣一層一層分叉的針葉樹）。 每次 \textsc{Draw} 在一條線的上方三分之二、左右兩側各長出
三條較短的斜線，再對其中一條線重複同樣的動作；同一個圖樣在越來越小的尺度上重複出現，是碎形（fractal）的畫法。
\end{pastsoln}
